% =====================================================================
%  Main-body paragraph + figures for ranking system performance.
%  Place after the Training section / before Verification Experiments.
%
%  Figures to upload to Overleaf  figures/  folder:
%    exp_bayesian_defended/figure1_diff_quality_6lines.png
%    exp_bayesian_defended/weight_trajectory.png
%    exp_bayesian_defended/figure2_similar_quality_planes_6lines.png
%    exp_bayesian_defended/figure4_quality_distribution.png
% =====================================================================

\section{Ranking Performance}
\label{sec:ranking-perf}

Figure~\ref{fig:rank-weight} (left) tracks the rank trajectories of six cat images spanning three quality tiers.
The two highest-quality images (Image~3, $q{=}0.725$; Image~7, $q{=}0.569$) converge to ranks~1 and~2 within the first 100~rounds and remain there for the duration of the simulation.
Mid-quality images settle into ranks~5--7, while the two lowest-quality images ($q{<}0.15$) drift below rank~8.
This monotonic separation demonstrates that the ranking system successfully \emph{infers} latent image quality solely from noisy user interactions---without ever observing the ground-truth quality score directly.
By round~500 the top-10 results for the ``cat'' query contain 10 out of 10 correctly categorised images, with the top-3 positions perfectly ordered by true quality.
Figure~\ref{fig:rank-weight} (right) shows that the learned policy weights stabilise after roughly 200~rounds: cosine similarity retains the largest share (${\approx}0.26$), while like rate, watch time, and social proof grow to balanced contributions (${\approx}0.20$ each), and freshness decays to its floor.
Additional figures---the quality distribution across categories and rank dynamics for similar-quality plane images---are provided in Appendix~\ref{app:extra-figs}.

\begin{figure}[htbp]
  \centering
  \begin{minipage}[t]{0.52\linewidth}
    \centering
    \includegraphics[width=\linewidth]{figures/figure1_diff_quality_6lines.png}
  \end{minipage}\hfill
  \begin{minipage}[t]{0.46\linewidth}
    \centering
    \includegraphics[width=\linewidth]{figures/weight_trajectory.png}
  \end{minipage}
  \caption{\textbf{Left:} Rank trajectories of six cat images at three quality tiers (high/mid/low).
           Higher-quality images converge to lower (better) ranks.
           \textbf{Right:} Policy weight trajectory over 1\,000~rounds.
           Weights stabilise after ${\sim}200$ rounds; cosine similarity dominates while engagement signals share the remaining budget.}
  \label{fig:rank-weight}
\end{figure}

% =====================================================================
%  Appendix figures (place in your appendix section)
% =====================================================================

% \appendix  % uncomment if not already in appendix mode
% \section{Additional Figures}
% \label{app:extra-figs}
%
% \begin{figure}[htbp]
%   \centering
%   \includegraphics[width=0.75\linewidth]{figures/figure4_quality_distribution.png}
%   \caption{Quality score distribution per category.  Each dot is one image;
%            horizontal lines mark the median.  Quality values are drawn
%            i.i.d.\ from $\mathrm{Beta}(2,5)$ and are hidden from the ranker.}
%   \label{fig:quality-dist}
% \end{figure}
%
% \begin{figure}[htbp]
%   \centering
%   \includegraphics[width=0.85\linewidth]{figures/figure2_similar_quality_planes_6lines.png}
%   \caption{Rank dynamics for six plane images with the tightest quality spread
%            ($q \in [0.075, 0.347]$).  Despite small quality differences, the
%            system gradually separates them, with higher-quality images
%            trending toward better ranks over 1\,000~rounds.}
%   \label{fig:planes-similar}
% \end{figure}
